{\large\textbf{E1 – Deep Learning (10 pts)}}

Modern ANNs (artificial neural networks) are made of billions of neurons. Each
neuron transforms its input(s) $x_1, x_2, \ldots, x_n$ to an output $y$. First,
\[
z = w_1x_1 + w_2x_2 + \cdots + w_nx_n + b
\]
is calculated, with real numbered \textit{weights} $w_i$ and real numbered
\textit{bias} $b$. Then an \textit{activation function} is applied to $z$ to
produce the final output $y(x_1, x_2, \ldots)$. In the present problem you will
investigate a physical model of a neuron with the electric voltages $x_1$ and
$x_2$ as inputs, with the activation function being $A\sigma(z)$, graphed
below, where $\sigma(z) = 1/(1{+}\exp\ (-z))$ is called sigmoid function.

\begin{center}\includegraphics[width=0.32\linewidth]{figures/EuPhO_2025_Q4_fig1.png}\end{center}

\textbf{Equipment}

\textbf{Important! Do not switch off any power outlets.}

(i)\quad A box containing a voltage source, an electronic circuit that models
the neuron, and two potentiometers (the A-potentiometer and the
B-potentiometer). The electric terminals on the box are denoted as follows:

1\quad Two electrically connected GND terminals: the electrical \textit{ground}
serving as a common negative terminal for +V, $x_1$, $x_2$, and $y$.

2\quad +V: the positive terminal of the voltage source.

3\quad X1 and X2: the positive terminals of the neuron input voltages $x_1$ and
$x_2$, respectively. \textbf{The neuron output behaves unpredictably if either
of these terminals has no input voltage.}

4\quad Y: the positive output terminal. It behaves like a real voltage source,
consisting of an ideal voltage source of voltage $y$ and a series output
resistor $R_\mathrm{out}$, and operates as shown below.

\begin{center}\includegraphics[width=0.54\linewidth]{figures/EuPhO_2025_Q4_fig2.png}\end{center}

5\quad A1, A2, A3: terminals of the A-potentiometer.

6\quad B1, B2, B3: terminals of the B-potentiometer.

7\quad T: a terminal not to be used in this task.

(ii)\quad Digital multimeter with two probe wires.

(iii)\quad Wires with banana connectors. Two or more wires could be connected
to the same terminal in the box by using the holes in the banana connectors.
Using the banana connectors with the multimeter may form an unstable
connection. Use the alligator clamp if needed.

(iv)\quad Graph paper. You can ask for more if needed.

\textbf{Task 1 (0.5 pts)}

Terminals A1, A2, and A3 are connected to the A-potentiometer $R_P$ and an
additional load resistor $R_L$.

Which of the schemes below corresponds to the circuit in the box? Determine
the resistances $R_L$ and $R_P$; document the measurements made.

\begin{center}\includegraphics[width=0.57\linewidth]{figures/EuPhO_2025_Q4_fig3.png}\end{center}

\textbf{Note}\quad The B-potentiometer is connected to terminals B1, B2, B3 in
exactly the same way with the same resistances $R_L$ and $R_P$, within
manufacturing tolerances.

\textbf{Task 2 - (0.5 pts)}

Sketch how the terminals have to be connected so that the neuron input voltages
can be varied with the widest possible range.

\textbf{Task 3 - (1.5 pts)}

Devise (and document) a strategy allowing you to find the combination of input
voltages $x_1$ and $x_2$ that maximizes the output voltage $y$ with the least
possible number of measurements, irrespectively of with which set of input
voltages you start the search. Determine this maximal voltage $y_\mathrm{max}$
that will be henceforth used as an approximation for the amplitude $A$, and
document your measurements.

\textbf{Task 4 - (3.5 pts)}

Determine the weights $w_1, w_2$ and the bias $b$. Describe your measurements
and document your data in a table. Estimate $w_1, w_2$, and $b$ by using a
graphical approach.

\textit{Training} involves optimizing the network weights to achieve desired
functionality. This allows ANNs to approximate arbitrary functions. For each of
the following tasks you have to approximate a different function of a single
input voltage using the given equipment. Make sure that the input and output
that you define are clearly marked in your circuits.

\textbf{Task 5 - (1.5 pts)}

Connect the terminal X1 directly to +V. Design a circuit to approximate the
function $y_5(x) = A\sigma\,(w_2x/2 + b_5)$, where $x$ is the voltage applied
to your newly defined input terminal. Determine $b_5$ theoretically. Implement
the circuit, take measurements and verify that your setup works as expected.
Validate the value of $b_5$ from your data.

\textbf{Task 6 - (2.5 pts)}

a\quad Determine the internal series output resistance $R_\mathrm{out}$ of the
Y terminal. (0.5 pts)

b\quad Design and implement a circuit to approximate the function
$y_6(x) = A_6 \cdot \sigma(w_2x + b) + B_6$, where
$B_6 = 1.48\,\mathrm{V}$. Determine $A_6$ theoretically. Implement the circuit
and verify experimentally that your setup works as expected. Confirm the
values of $A_6$ and $B_6$ from your data. (2.0 pts)
